% Generated by eval/scripts/write_bootstrap_example.py.
\paragraph*{A small example with actual scores.} To make the calculation concrete, consider Sagui-7B responses to just three source prompts, 0001, 0002, and 0003, in Portuguese and English. Their observed score pairs (Portuguese, English) are (0.875, 1.000), (0.250, 0.625), and (0.375, 0.250), respectively. Adding all six scores and dividing by six gives the original mean, $3.375/6 = 0.5625$.

If we randomly sample three source-prompt identifiers with replacement, one possible sample is 0001, 0001, and 0003. We include all recorded responses belonging to each selected identifier: the Portuguese and English scores for 0001 appear twice, those for 0003 appear once, and those for 0002 do not appear. This new sample still has six scores, but its mean is $4.375/6 \approx 0.7292$. It is higher because the draw included the higher-scoring prompt 0001 twice and omitted prompt 0002. Another draw could include prompt 0002 twice and give a lower mean. The scores in this example are real; the selected identifiers illustrate the procedure and are not reported analysis replicates.
